%!TEX program = lualatex
% Arquivo: relatorio_pratica_7_dureza_agua_mar.tex

\documentclass[a4paper,12pt]{article}

% =========================================================
% Idioma e tipografia
% =========================================================
\usepackage[brazilian]{babel}
\usepackage{fontspec}
\setmainfont{Latin Modern Roman}
\usepackage{microtype}

% =========================================================
% Layout
% =========================================================
\usepackage[top=3cm,bottom=2.5cm,left=3cm,right=2.5cm]{geometry}
\usepackage{setspace}
\onehalfspacing
\usepackage{indentfirst}
\setlength{\parindent}{1.25cm}

% =========================================================
% Matemática, química e unidades
% =========================================================
\usepackage{amsmath,amssymb}
\usepackage[version=4]{mhchem}
\usepackage{siunitx}
\sisetup{
  locale = DE,
  output-decimal-marker = {,},
  per-mode = symbol,
  exponent-product = \cdot
}

% =========================================================
% Tabelas, figuras e listas
% =========================================================
\usepackage{booktabs}
\usepackage{array}
\usepackage{graphicx}
\usepackage{float}
\usepackage{enumitem}
\usepackage{xcolor}
\usepackage[most]{tcolorbox}

\setlist[itemize]{noitemsep,topsep=3pt}
\setlist[enumerate]{noitemsep,topsep=3pt}

% =========================================================
% Links
% =========================================================
\usepackage{hyperref}
\hypersetup{
  colorlinks=true,
  linkcolor=black,
  urlcolor=blue,
  citecolor=black,
  pdfauthor={Israel Carneiro; José Mauro Xavier Elsas; Miguel Gayani},
  pdftitle={Prática 7 -- Determinação da dureza total da água do mar}
}

% =========================================================
% Dados experimentais
% =========================================================
\newcommand{\ConcEDTA}{0.009512}
\newcommand{\VolAmostra}{5.00}
\newcommand{\VolMedioTotal}{35.47}
\newcommand{\VolMedioCalcio}{6.20}

\begin{document}

% =========================================================
% Capa
% =========================================================
\begin{titlepage}
\centering

{\large Universidade do Estado do Rio de Janeiro -- UERJ\par}
\vspace{0.3cm}
{\large Química VIII / Oceanografia\par}

\vspace{1.2cm}

{\LARGE\bfseries Prática 7\par}
\vspace{0.4cm}
{\Large\bfseries DETERMINAÇÃO DA DUREZA TOTAL\\
DA ÁGUA DO MAR\par}

\vspace{0.8cm}

\IfFileExists{uerj.jpg}{
  \includegraphics[width=0.28\linewidth]{uerj.jpg}
}{
  \begin{tcolorbox}[
    width=0.30\textwidth,
    colback=white,
    colframe=black!35,
    boxrule=0.5pt,
    arc=0mm
  ]
  \centering
  \vspace{1.0cm}
  \textit{Logotipo da UERJ}
  \vspace{1.0cm}
  \end{tcolorbox}
}

\vfill

\begin{flushright}
\textbf{Alunos:}\\
Israel Carneiro\\
José Mauro Xavier Elsas\\
Miguel Gayani\\[0.5cm]
\textbf{Professor:} Ítalo Vieira
\end{flushright}

\vspace{1cm}
Rio de Janeiro -- 2026

\end{titlepage}

% =========================================================
\section{Introdução}

A dureza da água está relacionada principalmente à presença de cátions
metálicos multivalentes, com destaque para os íons cálcio, \ce{Ca^2+}, e
magnésio, \ce{Mg^2+}. Em água do mar, esses íons estão naturalmente
presentes em concentrações elevadas e participam de processos químicos e
biogeoquímicos importantes, como a formação de minerais carbonáticos e a
regulação do equilíbrio iônico do meio.

Outros cátions, como ferro, alumínio, cobre e zinco, também podem reagir
com o agente complexante e interferir na determinação. Caso estejam
presentes em quantidades relevantes, devem ser removidos ou mascarados antes
da análise. Neste experimento, a dureza total foi considerada como a soma
das contribuições de cálcio e magnésio, conforme o procedimento proposto na
apostila da disciplina.

A determinação da dureza pode ser realizada por volumetria de complexação
utilizando o ácido etilenodiaminotetracético (EDTA) como agente titulante.
O EDTA forma complexos estáveis com diversos íons metálicos. Para um cátion
divalente genérico, representado por \ce{M^2+}, a reação pode ser escrita,
de forma simplificada, como:

\begin{equation}
\ce{M^2+(aq) + Y^4-(aq) -> MY^2-(aq)},
\end{equation}

em que \ce{Y^4-} representa a forma do EDTA que participa da complexação.
A estequiometria da reação é de \textbf{1:1}; portanto, um mol de EDTA
complexa um mol de \ce{Ca^2+} ou um mol de \ce{Mg^2+}.

O EDTA atua como agente quelante hexadentado: seus dois átomos de nitrogênio
e quatro grupos carboxilato podem fornecer pares de elétrons para a formação
de ligações coordenadas com o íon metálico. O complexo resultante é
denominado quelato.

Os pontos finais das titulações foram identificados com indicadores
metalocrômicos. Esses indicadores formam complexos coloridos com os cátions,
mas o complexo metal--EDTA deve ser mais estável que o complexo
metal--indicador. Durante a titulação, o EDTA desloca o indicador:

\begin{equation}
\ce{M-Ind + EDTA -> M-EDTA + Ind}.
\end{equation}

As cores diferentes do indicador ligado e livre permitem reconhecer o ponto
final da titulação.

Neste experimento foram realizadas duas titulações. Na primeira,
determinou-se a quantidade total de \ce{Ca^2+} e \ce{Mg^2+} em meio
tamponado a pH 10, usando negro de eriocromo T como indicador. A viragem foi
observada pela mudança de vermelho para azul.

Na segunda titulação, determinou-se apenas o cálcio. A adição de solução de
\ce{NaOH} elevou o pH e retirou o magnésio da titulação por precipitação na
forma de \ce{Mg(OH)2}. Empregou-se calcon como indicador, com mudança de
rosa para azul. A concentração de magnésio foi então obtida pela diferença
entre a concentração total dos dois cátions e a concentração de cálcio.

% =========================================================
\section{Objetivos}

\begin{itemize}
  \item determinar a concentração total de \ce{Ca^2+} e \ce{Mg^2+} em uma
        amostra de água do mar por titulação complexométrica com EDTA;
  \item determinar separadamente a concentração de \ce{Ca^2+};
  \item calcular a concentração de \ce{Mg^2+} por diferença;
  \item expressar as concentrações em \si{\mole\per\litre}, percentual
        massa por volume (\% $m/V$) e ppm;
  \item expressar adicionalmente a dureza total em
        \si{\milli\gram\per\litre} como \ce{CaCO3}.
\end{itemize}

% =========================================================
\section{Material utilizado}

\begin{itemize}
  \item bureta de \qty{50.00}{\milli\litre};
  \item Erlenmeyers de \qty{250}{\milli\litre};
  \item pipeta aferida de \qty{5.00}{\milli\litre};
  \item pipeta graduada de \qty{5}{\milli\litre};
  \item proveta de \qty{50}{\milli\litre};
  \item béquer de \qty{400}{\milli\litre};
  \item bastão de vidro;
  \item suporte universal e garra para bureta;
  \item frasco lavador com água destilada.
\end{itemize}

% =========================================================
\section{Reagentes e soluções}

\begin{itemize}
  \item solução padrão de EDTA dissódico de concentração
        \qty{\ConcEDTA}{\mole\per\litre}, fornecida pronta pelo laboratório;
  \item amostra de água do mar;
  \item solução-tampão de pH 10;
  \item solução de \ce{NaOH} a 10\% ($m/V$);
  \item negro de eriocromo T;
  \item calcon;
  \item água destilada.
\end{itemize}

A apostila descreve o preparo da solução padrão a partir de EDTA dissódico
di-hidratado. Entretanto, nesta aula a solução foi fornecida pronta, com
concentração informada de \qty{\ConcEDTA}{\mole\per\litre}, pois a etapa de
preparo consumiria uma parcela excessiva do tempo disponível para a prática.

% =========================================================
\section{Procedimento experimental}

\subsection{Titulação conjunta de \ce{Ca^2+} e \ce{Mg^2+}}

Foram transferidos \qty{\VolAmostra}{\milli\litre} de água do mar para um
Erlenmeyer. Adicionaram-se \qty{50}{\milli\litre} de água destilada,
\qty{3}{\milli\litre} de solução-tampão de pH 10 e negro de eriocromo T.
A bureta foi rinsada e preenchida com a solução de EDTA. O Erlenmeyer foi
colocado sobre fundo branco e a amostra foi titulada com solução de EDTA
\qty{\ConcEDTA}{\mole\per\litre} até a mudança de coloração de vermelho
vinho para azul. Próximo ao ponto final, as paredes do Erlenmeyer foram
lavadas com água destilada e o titulante foi adicionado gota a gota. O
procedimento foi realizado em triplicata.

\subsection{Titulação de \ce{Ca^2+}}

Em outra alíquota de \qty{\VolAmostra}{\milli\litre} de água do mar,
adicionaram-se \qty{50}{\milli\litre} de água destilada,
\qty{2.5}{\milli\litre} de solução de \ce{NaOH} a 10\% ($m/V$) e calcon.
A solução foi agitada por dois minutos após a adição de \ce{NaOH}, para
favorecer a precipitação do magnésio como \ce{Mg(OH)2}. Em seguida,
adicionou-se o calcon e titulou-se com o mesmo EDTA, sobre fundo branco, até
a mudança de rosa para azul.

Depois de atingido o ponto final, a amostra permaneceu em repouso por cinco
minutos, com agitação ocasional. Essa etapa permite verificar se houve
formação de precipitado de \ce{Ca(OH)2}; caso a coloração vermelha
reaparecesse, a titulação deveria prosseguir até a obtenção de um novo ponto
final azul. A determinação foi realizada em triplicata.

Embora a relação geral de reagentes da apostila mencione solução de
\ce{NaOH} a 50\% ($m/V$), tanto o procedimento específico da Prática 7
quanto o quadro da aula indicavam \ce{NaOH} a 10\% ($m/V$). Por isso, foi
registrada no relatório a concentração efetivamente orientada para o
experimento.

A água destilada adicionada nas duas etapas apenas aumentou o volume da
solução no Erlenmeyer e facilitou a observação da viragem. Ela não alterou
a quantidade de matéria dos íons presentes na alíquota. Por isso, nos
cálculos das concentrações foi utilizado o volume original de
\qty{\VolAmostra}{\milli\litre} da amostra, e não o volume total após a
diluição no Erlenmeyer.

% =========================================================
\section{Resultados e cálculos}

\subsection{Volumes de EDTA consumidos}

\begin{table}[H]
\centering
\caption{Dados da titulação conjunta de \ce{Ca^2+} e \ce{Mg^2+}, em tabela
adaptada da apostila da disciplina.}
\label{tab:dados_total}
\begin{tabular}{lccc}
\toprule
\textbf{Dados da titulação} & \textbf{I} & \textbf{II} & \textbf{III}\\
\midrule
Concentração do EDTA (\si{\mole\per\litre})
  & 0,009512 & 0,009512 & 0,009512\\
Volume da amostra (\si{\milli\litre})
  & 5,00 & 5,00 & 5,00\\
Volume de titulante gasto (\si{\milli\litre})
  & 35,1 & 35,2 & 36,1\\
Ensaio em branco
  & \multicolumn{3}{c}{Não realizado}\\
Volume real de EDTA (\si{\milli\litre})
  & 35,1 & 35,2 & 36,1\\
\textbf{Volume médio real de EDTA (\si{\milli\litre})}
  & \multicolumn{3}{c}{\textbf{35,47}}\\
\bottomrule
\end{tabular}
\end{table}

\begin{table}[H]
\centering
\caption{Dados da titulação de \ce{Ca^2+}, em tabela adaptada da apostila
da disciplina.}
\label{tab:dados_calcio}
\begin{tabular}{lccc}
\toprule
\textbf{Dados da titulação} & \textbf{I} & \textbf{II} & \textbf{III}\\
\midrule
Concentração do EDTA (\si{\mole\per\litre})
  & 0,009512 & 0,009512 & 0,009512\\
Volume da amostra (\si{\milli\litre})
  & 5,00 & 5,00 & 5,00\\
Volume de titulante gasto (\si{\milli\litre})
  & 6,3 & 6,0 & 6,3\\
Ensaio em branco
  & \multicolumn{3}{c}{Não realizado}\\
Volume real de EDTA (\si{\milli\litre})
  & 6,3 & 6,0 & 6,3\\
\textbf{Volume médio real de EDTA (\si{\milli\litre})}
  & \multicolumn{3}{c}{\textbf{6,20}}\\
\bottomrule
\end{tabular}
\end{table}

Como não foi realizado ensaio em branco, não houve volume a subtrair dos
valores medidos. Desse modo, neste relatório:

\begin{equation}
V_{\text{real}}=V_{\text{gasto}}.
\end{equation}

Para a titulação conjunta:

\begin{align}
\overline{V}_{\ce{Ca + Mg}}
&=
\frac{35.1+35.2+36.1}{3}\\
&=
\qty{35.4667}{\milli\litre}
\approx
\boxed{\qty{35.47}{\milli\litre}}.
\end{align}

Para a titulação do cálcio:

\begin{align}
\overline{V}_{\ce{Ca}}
&=
\frac{6.3+6.0+6.3}{3}\\
&=
\boxed{\qty{6.20}{\milli\litre}}.
\end{align}

\subsection{Concentração total de \ce{Ca^2+} e \ce{Mg^2+}}

Como a proporção estequiométrica entre o EDTA e os cátions divalentes é
1:1:

\begin{equation}
n_{\mathrm{EDTA}}=n_{\ce{Ca^2+}}+n_{\ce{Mg^2+}}.
\end{equation}

Assim:

\begin{align}
[\ce{Ca^2+}]+[\ce{Mg^2+}]
&=
\frac{C_{\mathrm{EDTA}}\,\overline{V}_{\ce{Ca + Mg}}}
{V_{\mathrm{amostra}}}\\
&=
\frac{
\left(\qty{0.009512}{\mole\per\litre}\right)
\left(\qty{35.4667e-3}{\litre}\right)}
{\qty{5.00e-3}{\litre}}\\
&=
\boxed{\qty{0.06747}{\mole\per\litre}}.
\end{align}

\subsection{Concentração de \ce{Ca^2+}}

Na segunda titulação, o EDTA consumido corresponde apenas ao cálcio:

\begin{align}
[\ce{Ca^2+}]
&=
\frac{C_{\mathrm{EDTA}}\,\overline{V}_{\ce{Ca}}}
{V_{\mathrm{amostra}}}\\
&=
\frac{
\left(\qty{0.009512}{\mole\per\litre}\right)
\left(\qty{6.20e-3}{\litre}\right)}
{\qty{5.00e-3}{\litre}}\\
&=
\boxed{\qty{0.01179}{\mole\per\litre}}.
\end{align}

\subsection{Concentração de \ce{Mg^2+}}

A concentração de magnésio foi calculada por diferença:

\begin{align}
[\ce{Mg^2+}]
&=
\left([\ce{Ca^2+}]+[\ce{Mg^2+}]\right)-[\ce{Ca^2+}]\\
&=
0.0674718-0.0117949\\
&=
\boxed{\qty{0.05568}{\mole\per\litre}}.
\end{align}

% =========================================================
\subsection{Conversão para percentual massa por volume}

O percentual massa por volume indica a massa, em gramas, do soluto contida
em \qty{100}{\milli\litre} de solução. Para uma concentração molar $C$ e
massa molar $M$:

\begin{equation}
\%\,(m/V)=\frac{C\left(\si{\mole\per\litre}\right)
M\left(\si{\gram\per\mole}\right)}{10}.
\end{equation}

Foram utilizadas as massas molares
$M_{\ce{Ca}}=\qty{40.08}{\gram\per\mole}$ e
$M_{\ce{Mg}}=\qty{24.31}{\gram\per\mole}$.

Para o cálcio:

\begin{align}
\%\,(m/V)_{\ce{Ca^2+}}
&=
\frac{(0.0117949)(40.08)}{10}\\
&=
\boxed{0.04727\%\ (m/V)}.
\end{align}

Para o magnésio:

\begin{align}
\%\,(m/V)_{\ce{Mg^2+}}
&=
\frac{(0.0556769)(24.31)}{10}\\
&=
\boxed{0.1354\%\ (m/V)}.
\end{align}

% =========================================================
\subsection{Conversão para ppm}

Para soluções aquosas diluídas, a concentração em ppm pode ser expressa
numericamente em \si{\milli\gram\per\litre}. A conversão foi realizada por:

\begin{equation}
C(\mathrm{ppm})=
C\left(\si{\mole\per\litre}\right)
M\left(\si{\gram\per\mole}\right)
\times 1000.
\end{equation}

Para o cálcio:

\begin{align}
C_{\ce{Ca^2+}}
&=
(0.0117949)(40.08)(1000)\\
&=
\boxed{\qty{472.7}{\milli\gram\per\litre}}
\approx
\boxed{472.7\ \mathrm{ppm}}.
\end{align}

Para o magnésio:

\begin{align}
C_{\ce{Mg^2+}}
&=
(0.0556769)(24.31)(1000)\\
&=
\boxed{\qty{1354}{\milli\gram\per\litre}}
\approx
\boxed{1354\ \mathrm{ppm}}.
\end{align}

% =========================================================
\subsection{Dureza total expressa como \ce{CaCO3}}

Embora o quadro solicite as concentrações individuais de cálcio e magnésio,
a dureza também é convencionalmente expressa em
\si{\milli\gram\per\litre} como \ce{CaCO3}. Como um mol de cada cátion
divalente corresponde, em termos de equivalência, a um mol de \ce{CaCO3}:

\begin{align}
\text{Dureza total}
&=
\left([\ce{Ca^2+}]+[\ce{Mg^2+}]\right)
M_{\ce{CaCO3}}\times1000\\
&=
(0.0674718)(100.09)(1000)\\
&=
\boxed{\qty{6753}{\milli\gram\per\litre}
\ \text{como \ce{CaCO3}}}.
\end{align}

% =========================================================
\subsection{Resumo dos resultados}

\begin{table}[H]
\centering
\caption{Concentrações de cálcio e magnésio na amostra de água do mar.}
\label{tab:resultados}
\begin{tabular}{
l
S[table-format=1.5]
S[table-format=1.5]
S[table-format=4.1]
}
\toprule
\textbf{Íon}
& {\textbf{Concentração} (\si{\mole\per\litre})}
& {\textbf{\%} ($m/V$)}
& {\textbf{ppm}}\\
\midrule
\ce{Ca^2+} & 0,01179 & 0,04727 & 472,7\\
\ce{Mg^2+} & 0,05568 & 0,1354  & 1354,0\\
\bottomrule
\end{tabular}
\end{table}

% =========================================================
\section{Discussão}

A titulação conjunta consumiu, em média,
\qty{35.47}{\milli\litre} de EDTA, enquanto a titulação seletiva do cálcio
consumiu \qty{6.20}{\milli\litre}. Como o EDTA reage na proporção 1:1 com
os dois cátions, a diferença entre esses volumes está associada ao magnésio.
Consequentemente, a concentração molar de \ce{Mg^2+} foi superior à de
\ce{Ca^2+} na amostra analisada.

As concentrações calculadas foram de
\qty{0.01179}{\mole\per\litre} para o cálcio e
\qty{0.05568}{\mole\per\litre} para o magnésio. Em concentração mássica,
esses valores correspondem a aproximadamente \qty{472.7}{\milli\gram\per\litre}
de cálcio e \qty{1354}{\milli\gram\per\litre} de magnésio. Os resultados
apresentam a predominância do magnésio característica da composição iônica
da água do mar.

Na titulação conjunta, o tampão de pH 10 forneceu as condições adequadas
para a complexação dos dois cátions pelo EDTA e para o funcionamento do
negro de eriocromo T. Na titulação seletiva, a adição de \ce{NaOH} tornou
o meio fortemente básico, promovendo a precipitação do magnésio como
\ce{Mg(OH)2}. Dessa forma, o EDTA permaneceu disponível para reagir
predominantemente com o cálcio, e a viragem foi acompanhada com calcon.

Os volumes da primeira titulação variaram entre
\qty{35.1}{\milli\litre} e \qty{36.1}{\milli\litre}, enquanto os da
segunda variaram entre \qty{6.0}{\milli\litre} e
\qty{6.3}{\milli\litre}. Apesar dessas diferenças, os três resultados de
cada etapa foram mantidos no cálculo das médias, pois não havia evidência
experimental suficiente para classificar qualquer medida como resultado
anômalo.

A diluição com \qty{50}{\milli\litre} de água destilada não foi incluída
no denominador das equações de concentração. Isso ocorre porque a diluição
não modifica a quantidade de matéria de cálcio e magnésio originalmente
presente nos \qty{5.00}{\milli\litre} da amostra; modifica apenas o volume
ocupado pela solução durante a titulação.

A solução de EDTA não foi preparada durante a aula. Ela foi fornecida pronta
pelo laboratório, com concentração de
\qty{0.009512}{\mole\per\litre}, valor informado no quadro para as duas
titulações. Essa adaptação permitiu dedicar o tempo da aula às determinações
complexométricas e não modifica os cálculos, desde que seja empregada a
concentração efetiva informada para o titulante.

O roteiro prevê a determinação de um volume de branco, que deveria ser
subtraído do volume gasto para obtenção do volume real de EDTA. Como o ensaio
em branco não foi realizado nesta aula, adotou-se
$V_{\text{real}}=V_{\text{gasto}}$. A ausência dessa correção constitui uma
limitação do experimento: eventual consumo de titulante por reagentes ou por
outras contribuições do procedimento permanece incorporado aos volumes
medidos e pode produzir um pequeno erro sistemático positivo nas
concentrações calculadas.

% =========================================================
\section{Conclusão}

A dureza total da amostra de água do mar foi determinada por titulação
complexométrica com solução de EDTA
\qty{0.009512}{\mole\per\litre}. A primeira titulação forneceu a soma das
concentrações de cálcio e magnésio:

\begin{equation*}
\boxed{
[\ce{Ca^2+}]+[\ce{Mg^2+}]
=
\qty{0.06747}{\mole\per\litre}
}.
\end{equation*}

A titulação seletiva do cálcio e o cálculo do magnésio por diferença
produziram os seguintes resultados:

\begin{align*}
[\ce{Ca^2+}]
&=
\boxed{\qty{0.01179}{\mole\per\litre}}
=
\boxed{472.7\ \mathrm{ppm}},\\[0.2cm]
[\ce{Mg^2+}]
&=
\boxed{\qty{0.05568}{\mole\per\litre}}
=
\boxed{1354\ \mathrm{ppm}}.
\end{align*}

Em percentual massa por volume, foram obtidos 0,04727\% ($m/V$) de cálcio
e 0,1354\% ($m/V$) de magnésio. A dureza total, expressa de forma
convencional como \ce{CaCO3}, foi de aproximadamente
\qty{6753}{\milli\gram\per\litre}. Os resultados mostraram concentração de
magnésio superior à de cálcio, comportamento compatível com a composição
iônica esperada para a água do mar. Como não foi realizado ensaio em branco,
os volumes reais foram considerados iguais aos volumes gastos; portanto, os
resultados devem ser interpretados levando em conta essa limitação
experimental.

% =========================================================
\section*{Referência}

LAGO, Dalva Cristina Baptista do; GOIS, Jefferson Santos de; LANGONE,
Marta Antunes Pereira; SENNA, Lilian Ferreira de.
\textit{Química VIII -- Oceanografia: aulas práticas}.
Rio de Janeiro: Universidade do Estado do Rio de Janeiro,
Instituto de Química, Departamento de Química Analítica, 2017.

\end{document}
